\paragraph{Replica prescriptions against published contact-term treatments (conditional; not closure).}
Keep the setting, the fixed-state replica and the first-order smooth replica of the preceding paragraph, and write \(c_A\) and \(c_A^{\mathrm{sm}}\) for their area coefficients and \(c_R=N_{\mathrm{tot}}/6\) for the coefficient of \(\rho R[g_\rho]\); \(c_A=2\) and \(c_A^{\mathrm{sm}}=0\) at \(n_D=2\), \(N_s=6\). Put \(V_\mu=\partial_\mu\ln\rho\).

\emph{Scalar and ghost sectors.} For a scalar with a coupling \(\xi R\phi^2\), Herzog and Nishioka add the codimension-two term \(\delta I_E=-2\pi(1-n)\xi\int\phi^2\) (Eq.~(2.19) of~\cite{HerzogNishioka2016}), which cancels the delta of \(\xi R\) on the cone and gives an entropy independent of \(\xi\). With the heat-trace bookkeeping of the contact-term paragraph above, this term shifts the entropy by \(+2\pi\xi A\langle\phi^2\rangle\), equal and opposite to the contact \(-2\pi\xi A\langle\phi^2\rangle\), which is the expectation value of Wald's formula for \(\xi R\phi^2\)~\cite{DonnellyWall2016}, Eq.~(4); Herzog and Nishioka do not discuss Wald's formula, and their term is related to it only through this bookkeeping. The fixed-state replica, whose scalar and ghost operators carry no delta on the cone, therefore has the contact-free coefficient. The smooth replica, in which \(E_s\) contains \(-\tfrac12\Box_G\ln\rho\) with \(Q=\tfrac14\), has the form of the delta-potential entropy of Eqs.~(143)--(144) of~\cite{Solodukhin2011} with \(\xi\) replaced by \(Q\); that review does not accept this entropy as an entanglement entropy, and Donnelly and Wall assign the corresponding term to the generalised entropy as a Wald contact. Neither value is asserted here to be the entanglement entropy of the matrix state: the identifications are by the form of the coefficient, the cited analysis is for a scalar on a flat background, and the operators here have no independent \(\xi\). On the singular-\(\rho\) cone \(\ln\rho=(n-1)\ln|z|\), the operator \(\mathcal O_s\) is unitarily equivalent to \(-\Box+E_s\), with s-wave Bessel index \(\nu_0=|n-1|/(2n)>0\). The regular branch vanishes at \(\Sigma\), so the delta multiplies a function that vanishes there, and the constant \(-\nu_0/2\) gives the entropy coefficient \(\tfrac16-\tfrac14\) per component for \(n<1\) and \(\tfrac16+\tfrac14\) for \(n>1\), non-analytic at \(n=1\); the less singular branch exchanges the two. The smooth value \(-\tfrac1{12}\) is thus the \(n<1\) value of the regular branch, and removing the delta alone does not carry this cone to the contact-free value \(\tfrac16\). The terms of order \(((n-1)/2)^2\) in the indices \(\nu_k\) with \(k\neq0\) depend logarithmically on the tip scale and drop out at first order in \(n-1\).

\emph{Connection sector.} For the Maxwell field in Feynman gauge, Herzog and Nishioka remove the Ricci delta of the vector operator with \(\delta I=-\pi(1-n)\int\delta^{ij}A_iA_j\) (Eq.~(7.2) of~\cite{HerzogNishioka2016}), and note that this term is not gauge invariant. Kabat's contact is that delta: it is present in the black-hole entropy and absent from the bulk-mode entanglement entropy~\cite{Kabat1995}, and Donnelly and Wall identify it with edge modes~\cite{DonnellyWall2016}. Herzog and Nishioka state that their boundary term contributes oppositely to the contact term (Eq.~(5.20) of~\cite{HerzogNishioka2016}). The vector contact \(-1\) that gives \(c_A=1\) is therefore the term that the Herzog--Nishioka term removes. Its use presupposes a \(U(1)\) structure and a Ricci coupling that the operator \(\mathcal O_v\) does not have in the variables \((G,\rho,\partial\Theta)\), where the gauge group is symplectic and \(Q=0\); it remains an input and is not a prescription of this class.

\emph{Teleparallel form.} The identity for the Weitzenb\"ock connection of the paragraph on the torsion invariant gives \(E_s=-\tfrac12\nabla_\mu V^\mu+\tfrac14V_\mu V^\mu=\tfrac14\bigl(R_G+\mathcal T\bigr)+\tfrac14V_\mu V^\mu\) with \(\mathcal T=\tfrac14T_{\lambda\mu\nu}T^{\lambda\mu\nu}+\tfrac12T_{\lambda\mu\nu}T^{\mu\lambda\nu}-V_\mu V^\mu\), and the scalar density is \(\tfrac16R_G-E_s=-\tfrac1{12}R_G-\tfrac14\mathcal T-\tfrac14V_\mu V^\mu\). With \(\mathcal T=-R_G-2\nabla_\mu V^\mu\) the coefficient of \(R_G\) is \(\tfrac16\), which is the value in the quadratic form on the two modes with \(\delta\rho=0\), where \(V_\mu V^\mu\) is of higher order. For scalar content alone \(c_T=c_A\), as stated above, and the two replicas differ by the integrated boundary term \(R_G+\mathcal T=-2\nabla_\mu V^\mu\), the flux of \(\nabla\ln\rho\) through \(\Sigma\).

\emph{Dynamical metric.} In the replica construction of Lewkowycz and Maldacena the entropy is \(-\partial_n\) of the action of a regularised cone minus that of the cone, localised at the tip; the difference between the action of the true \(n\)-solution and that of the regularised cone is a first-order variation about a solution and vanishes by the equations of motion (Eqs.~(3.1)--(3.4) of~\cite{LewkowyczMaldacena2013}). That step requires the background to solve the equations of motion of the full action, of which the stationarity condition of the tensor block is the part not derived here. The regularised cone has the far field of the \(n=1\) solution with the angle identified modulo \(2\pi n\), so that every field other than the metric is regular there. With \(\Theta\) the only field, \(R_G^{(1)}=-2\Box_{\bar G}\,\delta\ln\rho\) excludes such a configuration at first order, as stated above. The only first-order regularisation in this class is the smooth family, whose far field has \(\rho\propto r^{n-1}\) and so differs from that of the cone; the construction of~\cite{LewkowyczMaldacena2013}, and the vanishing of its first bracket, are therefore not reproduced for it. For the smooth replica the area coefficient is \(c_A^{\mathrm{sm}}=-c_R/2\). This is the coefficient of \(\sqrt G\,R_G\) at fixed coframe torsion, so the smooth replica satisfies \(S=A/4G_{\mathrm{ind}}\) in those variables for every content, whereas \(c_A^{\mathrm{sm}}-c_T=(14-N_s-4n_D)/4\), which vanishes only for \(N_s+4n_D=14\). The generalised-entropy structure therefore does not imply \(G_A=G_T\) in this class. At \(n_D=2\), \(N_s=6\), the conditions \(c_R=0\) (\(N_s=4n_D-2\)) and \(c_T=0\) (\(N_s=22-8n_D\)) are independent and meet only there; \(c_A^{\mathrm{sm}}=0=c_T\) is then the absence of the \(\LUV^2\) term in the induced action and carries no information on \(G\), while \(c_A=2\) is the contact-free entanglement coefficient, in the relation, in form, \(S_{\mathrm{BH}}=S_{\mathrm{ent}}-Q\) of Eq.~(322) of~\cite{Solodukhin2011}, where \(Q\) is here the bosonic contact \((N_s+2)/4\) of the smooth reading and \(S_{\mathrm{BH}}\) is built from the induced action.

\emph{Intermediate flux.} For first-order smooth perturbations of constant \(\bar\Theta\) the integral of \(R_G^{(1)}\) over the transverse plane is \(-2\) times the flux of \(\nabla\ln\rho\), so no family of such perturbations joins the two replicas. If the first-order local densities of the pulled-back cone and of the smooth family are superposed with weight \(f\) on the smooth family, then \(c_A(f)=c_A-f(N_s+2)/4\), which is \(2-2f\) at \(n_D=2\), \(N_s=6\); the superposition is an assumption. The indices \(\nu_k=\sqrt{k^2+((n-1)/2)^2}/n=|k+\mathrm i(n-1)/2|/n\) are those of an imaginary twist, so the periodic quadratic dependence of the area coefficient on a real Abelian holonomy is not used.

These statements are for the Poisson-limit operators, the \(\LUV^2\) term, one loop, \(\alpha=1\), and the two replicas stated; they do not derive the replica or the state selected by the matrix model, the stationarity condition, the covariance of the induced action in \(\Theta\), the vector contact, or the tip domain, they treat neither edge modes of the symplectic gauge sector nor fermion edge modes, and they give no value or sign of \(G_{\mathrm{eff}}\). Problems~B, C, D, F and~G remain open.
